% !TEX program = xelatex
\documentclass[11pt,a4paper]{article}
\usepackage{algo-preamble}

\title{\textbf{L08 \textperiodcentered{} Γραφήματα III}\\[2pt]
\large Εφαρμογές, Διμερότητα, Κατευθυνόμενα}
\author{Algorithms Class Hub}
\date{\small Αλγόριθμοι και Πολυπλοκότητα (K17) \textperiodcentered{} ΕΚΠΑ}

\begin{document}
\maketitle

\begin{center}\itshape\small
Εφαρμογές διάσχισης: συνεκτικές συνιστώσες, flood-fill, έλεγχος διμερότητας με
BFS, κατευθυνόμενα γραφήματα και ισχυρή συνεκτικότητα. Εισαγωγή στο πρόβλημα
συντομότερης διαδρομής με βάρη.
\end{center}

\section{Τι θα δούμε}

Έχουμε BFS και DFS από το προηγούμενο μάθημα. Σε αυτή τη διάλεξη βλέπουμε
\textbf{πώς τα χρησιμοποιούμε} για να λύσουμε πραγματικά προβλήματα:
\begin{enumerate}
  \item \textbf{Συνεκτικές συνιστώσες} --- βρες όλες, χωρίς να ξέρεις από πού να
  ξεκινήσεις.
  \item \textbf{Flood fill} --- γέμισμα περιοχών με χρώμα σε εικόνα.
  \item \textbf{Έλεγχος διμερότητας} --- μπορώ να χρωματίσω το γράφημα με 2
  χρώματα ώστε κάθε ακμή να έχει διαφορετικά χρώματα στα δύο της άκρα;
  \item \textbf{Κατευθυνόμενα γραφήματα} --- γραφήματα όπου οι ακμές έχουν
  κατεύθυνση. Παγκόσμιος Ιστός, χρονοπρογραμματισμός, εξάρτηση εργασιών.
  \item \textbf{Ισχυρή συνεκτικότητα} σε κατευθυνόμενα γραφήματα.
  \item \textbf{Εισαγωγή στο πρόβλημα συντομότερης διαδρομής με βάρη} ---
  προετοιμασία για το επόμενο μάθημα όπου θα δούμε Dijkstra και Bellman-Ford.
\end{enumerate}

\section{Εφαρμογή 1: Συνεκτικές συνιστώσες}

\textbf{Πρόβλημα.} Δεδομένου γραφήματος $G = (V, E)$, βρες \textbf{όλες} τις
συνεκτικές συνιστώσες --- δηλαδή τις μέγιστες ομάδες κορυφών όπου κάθε δύο
συνδέονται με διαδρομή.

\begin{codebox}
\begin{verbatim}
ConnectedComponents(G):
  Μαρκάρισε όλες τις κορυφές ως «μη επισκεφθείσες»
  components <- []
  Για κάθε v ∈ V:
    Εάν v «μη επισκεφθείσα»:
      C = BFS(G, v)             // ή DFS(G, v)
      components.append(C)
  επίστρεψε components
\end{verbatim}
\end{codebox}

\textbf{Παράδειγμα:} ένα γράφημα 13 κορυφών μπορεί να χωρίζεται σε τρεις
συνιστώσες --- π.χ.\ $\{1, \dots, 8\}$, $\{9, 10\}$ και $\{11, 12, 13\}$.

\textbf{Πολυπλοκότητα:} συνολικά $O(n + m)$. Κάθε κορυφή και ακμή εξετάζεται
\textbf{μία φορά} σε όλο τον αλγόριθμο --- ακόμα κι αν το εξωτερικό loop τρέξει
$n$ φορές, οι περισσότερες κορυφές ήδη έχουν επισκεφθεί.

\section{Εφαρμογή 2: Flood fill (γέμισμα περιοχών με χρώμα)}

\textbf{Πρόβλημα.} Σε μια εικόνα έχουμε πράσινα εικονοστοιχεία (pixels). Δίνεται
ένα συγκεκριμένο πράσινο pixel --- αλλάξε \textbf{όλη} την επιφάνεια
\textbf{συνεκτικά γειτονικών} πράσινων pixels σε μπλε.

\textbf{Μοντελοποίηση σε γράφημα:}
\begin{itemize}
  \item \textbf{Κορυφές:} τα pixels της εικόνας.
  \item \textbf{Ακμή} $\{p, q\}$: τα $p, q$ είναι γειτονικά
  (πάνω/κάτω/δεξιά/αριστερά) \textbf{και} τα δύο πράσινα.
  \item \textbf{Επιφάνεια:} η συνεκτική συνιστώσα.
\end{itemize}
\textbf{Λύση:} BFS ή DFS από το δοσμένο pixel, χρωμάτισε ό,τι επισκεφθείς.
$O(n + m)$ όπου $n$ το πλήθος pixels και $m$ το πλήθος γειτνιακών ζευγών.

\begin{intuition}
\textbf{Στο Photoshop / MS Paint} το «paint bucket» tool είναι ακριβώς αυτό.
Κάθε φορά που πατάς paint bucket τρέχει BFS από το pixel που πάτησες.
\end{intuition}

\section{Εφαρμογή 3: Έλεγχος διμερότητας}

\subsection{Τι σημαίνει διμερές}

\textbf{Διμερές (bipartite)} γράφημα: μπορείς να χρωματίσεις τις κορυφές με 2
χρώματα (π.χ.\ \textbf{κόκκινο} και \textbf{μπλε}) ώστε κάθε ακμή να έχει
\textbf{διαφορετικό χρώμα} στα δύο της άκρα.

\textbf{Διαισθητικά:} οι κορυφές χωρίζονται σε \textbf{δύο σύνολα} $L$ και $R$,
και όλες οι ακμές πάνε από το $L$ στο $R$ (καμία μέσα στο $L$ ή μέσα στο $R$).

\subsection{Γιατί μας νοιάζει}

\begin{itemize}
  \item \textbf{Ευσταθές ταίριασμα} (γιατροί $\leftrightarrow$ νοσοκομεία):
  κόκκινες κορυφές $=$ νοσοκομεία, μπλε $=$ γιατροί, ακμή $=$ δυνατότητα
  τοποθέτησης.
  \item \textbf{Χρονοπρογραμματισμός} (μηχανές $\leftrightarrow$ εργασίες).
  \item \textbf{Πολλά προβλήματα γραφημάτων} που είναι NP-δύσκολα στο γενικό
  γράφημα γίνονται \textbf{πολυωνυμικά} στα διμερή (π.χ.\ μέγιστο ταίριασμα ---
  λύνεται με flow networks).
\end{itemize}

\begin{intuition}
Σκέψου πάντα: «μπορώ να χωρίσω τους κόμβους σε δύο ομάδες ώστε οι σχέσεις να
είναι μόνο μεταξύ ομάδων;» Αν ναι, το πρόβλημά σου ίσως είναι ευκολότερο από όσο
φαίνεται.
\end{intuition}

\subsection{Εμπόδιο: περιττοί κύκλοι}

\begin{keypoint}
\textbf{Λήμμα.} Αν γράφημα $G$ είναι διμερές, δεν μπορεί να περιέχει \textbf{κύκλο
περιττού μήκους}.

\textbf{Απόδειξη.} Σε έναν κύκλο περιττού μήκους $k$, αν χρωματίσουμε τις κορυφές
με 2 χρώματα κυκλικά, στο τέλος καταλήγουμε σε αντίφαση: η τελευταία ακμή ενώνει
δύο κορυφές ίδιου χρώματος.
\end{keypoint}

\subsection{Το αντίστροφο επίσης ισχύει}

\begin{keypoint}
\textbf{Θεώρημα (κύριο).} Έστω $G$ συνεκτικό γράφημα. Τρέξε BFS από οποιαδήποτε
κορυφή $s$ και πάρε επίπεδα $L_0, L_1, \dots, L_k$. \textbf{Ακριβώς ένα} από τα
παρακάτω ισχύει:
\begin{enumerate}
  \item Καμία ακμή του $G$ δεν ενώνει κορυφές του ίδιου επιπέδου $\Rightarrow$
  \textbf{το $G$ είναι διμερές} (χρωμάτισε άρτια επίπεδα κόκκινα, περιττά μπλε).
  \item Υπάρχει ακμή του $G$ που ενώνει δύο κορυφές του ίδιου επιπέδου
  $\Rightarrow$ \textbf{το $G$ έχει κύκλο περιττού μήκους}, άρα \textbf{όχι}
  διμερές.
\end{enumerate}
\end{keypoint}

\subsection{Απόδειξη του θεωρήματος}

\textbf{Περίπτωση 1:} αν καμία ακμή δεν είναι μέσα σε ίδιο επίπεδο, τότε από την
ιδιότητα του BFS-δέντρου που είδαμε στο προηγούμενο μάθημα («τα επίπεδα των άκρων
μιας ακμής διαφέρουν το πολύ κατά 1») κάθε ακμή πάει από επίπεδο $i$ σε επίπεδο
$i+1$. Χρωματίζοντας τα άρτια επίπεδα \textbf{κόκκινα} και τα περιττά
\textbf{μπλε}, \textbf{κάθε ακμή} έχει διαφορετικά χρώματα στα άκρα της.
\textbf{Διμερές.} $\checkmark$

\textbf{Περίπτωση 2:} έστω ακμή $(x, y)$ με $x, y$ στο ίδιο επίπεδο $L_j$. Έστω
$z = \text{LCA}(x, y)$ ο \textbf{χαμηλότερος κοινός πρόγονος} των $x, y$ στο
BFS-δέντρο, στο επίπεδο $L_i$ (με $i < j$). Σχηματίζουμε κύκλο:
\begin{itemize}
  \item την ακμή $(x, y)$ --- μήκος \textbf{1},
  \item το μονοπάτι $y \to z$ --- μήκος $j - i$,
  \item το μονοπάτι $z \to x$ --- μήκος $j - i$.
\end{itemize}
Συνολικό μήκος: $1 + (j - i) + (j - i) = 1 + 2(j - i)$ --- \textbf{περιττό}. Άρα
έχουμε περιττό κύκλο $\Rightarrow$ από το προηγούμενο λήμμα, \textbf{όχι}
διμερές. $\blacksquare$

\begin{keypoint}
\textbf{Συνέπεια --- αλγόριθμος διμερότητας:} τρέξε BFS, και κατά τη σάρωση κάθε
ακμής έλεγξε αν τα άκρα είναι στο ίδιο επίπεδο. Αν ναι, \textbf{όχι} διμερές.
Αλλιώς, διμερές. \textbf{$O(n + m)$.}
\end{keypoint}

\section{Κατευθυνόμενα Γραφήματα}

\subsection{Ορισμός}

Στα κατευθυνόμενα γραφήματα οι ακμές έχουν \textbf{κατεύθυνση}: $G = (V, A)$ όπου
$A$ είναι \textbf{διατεταγμένα ζευγάρια} $(u, v)$ --- η ακμή κατευθύνεται
\textbf{από} $u$ \textbf{προς} $v$.

\textbf{Νέοι ορισμοί βαθμών:}
\begin{itemize}
  \item \textbf{Εσώβαθμος} $\text{indeg}(v)$: πλήθος ακμών που μπαίνουν στην $v$.
  \item \textbf{Εξώβαθμος} $\text{outdeg}(v)$: πλήθος ακμών που βγαίνουν από την
  $v$.
\end{itemize}

\textbf{Λήμμα:} $\sum_v \text{indeg}(v) = \sum_v \text{outdeg}(v) = |A|$. Κάθε
ακμή μετράει μία φορά στους εσώβαθμους και μία στους εξώβαθμους.

\subsection{Παραδείγματα}

\begin{itemize}
  \item \textbf{Παγκόσμιος Ιστός:} κορυφές $=$ ιστοσελίδες, ακμές $=$
  υπερσύνδεσμοι. Η κατεύθυνση \textbf{έχει σημασία} --- μια σελίδα μπορεί να έχει
  link προς άλλη χωρίς να ισχύει το αντίστροφο.
  \item \textbf{Χρονοπρογραμματισμός εργασιών:} ακμή $a \to b$ σημαίνει «η $a$
  πρέπει να γίνει \textbf{πριν} την $b$».
  \item \textbf{Τροφική αλυσίδα:} $a \to b$ σημαίνει «το $a$ τρώει το $b$».
  \item \textbf{Καλέσματα συναρτήσεων} ενός προγράμματος: $f \to g$ σημαίνει «η
  $f$ καλεί τη $g$».
\end{itemize}

\subsection{Διάσχιση κατευθυνόμενου γραφήματος}

Ο \textbf{BFS} (και ο DFS) επεκτείνεται φυσικά σε κατευθυνόμενα γραφήματα: όταν
εξετάζουμε την κορυφή $u$, ακολουθούμε \textbf{μόνο} τις ακμές \textbf{που
βγαίνουν} από την $u$ --- δηλαδή τους εξώβαθμους γείτονες. Η πολυπλοκότητα
παραμένει $O(n + m)$ όπου $m = |A|$ το πλήθος κατευθυνόμενων ακμών.

\textbf{Πρόβλημα:} κατευθυνόμενη προσπελασιμότητα --- δεδομένης κορυφής $s$, βρες
όλες τις κορυφές που είναι προσπελάσιμες \textbf{ακολουθώντας} τις κατευθύνσεις.

\textbf{Εφαρμογή: Web crawler} --- ξεκίνα από την ιστοσελίδα $s$, ακολούθα όλα τα
links, βρες όλες τις προσπελάσιμες σελίδες.

\section{Ισχυρή Συνεκτικότητα}

\subsection{Ορισμός}

Σε κατευθυνόμενο γράφημα, δύο κορυφές $u, v$ είναι \textbf{αμοιβαία
προσπελάσιμες} αν υπάρχει διαδρομή $u \to v$ \textbf{και} διαδρομή $v \to u$ (όχι
κατ' ανάγκη διακεκριμένες μεταξύ τους --- μπορεί να μοιράζονται κορυφές).

Ένα κατευθυνόμενο γράφημα είναι \textbf{ισχυρά συνεκτικό} (strongly connected)
αν \textbf{κάθε ζευγάρι} κορυφών είναι αμοιβαία προσπελάσιμο.

\subsection{Πώς το ελέγχουμε σε γραμμικό χρόνο}

\begin{keypoint}
\textbf{Λήμμα.} Σε κατευθυνόμενο γράφημα $G$, διάλεξε \textbf{οποιαδήποτε} κορυφή
$s$. Το $G$ είναι ισχυρά συνεκτικό \textbf{αν και μόνο αν}:
\begin{enumerate}
  \item Κάθε κορυφή είναι προσπελάσιμη από την $s$ (στο $G$).
  \item Η $s$ είναι προσπελάσιμη από κάθε κορυφή.
\end{enumerate}

\textbf{Απόδειξη.} \textbf{($\Rightarrow$)} Απευθείας από τον ορισμό --- αν είναι
ισχυρά συνεκτικό, οποιοδήποτε ζεύγος είναι αμοιβαία προσπελάσιμο, άρα κάθε κορυφή
$\leftrightarrow$ $s$.

\textbf{($\Leftarrow$)} Πάρε οποιεσδήποτε $u, v$. Από συνθήκη 1, υπάρχει διαδρομή
$s \to v$. Από συνθήκη 2, υπάρχει διαδρομή $u \to s$. Συνένωση: $u \to s \to v$
--- διαδρομή από $u$ σε $v$. Συμμετρικά: $v \to s \to u$. Άρα αμοιβαία
προσπελάσιμοι. $\blacksquare$
\end{keypoint}

\subsection{Αλγόριθμος}

\textbf{Θεώρημα.} Μπορούμε να αποφασίσουμε αν ένα κατευθυνόμενο $G$ είναι ισχυρά
συνεκτικό σε \textbf{χρόνο $O(n + m)$}.

\begin{codebox}
\begin{verbatim}
StronglyConnected(G):
  s <- οποιαδήποτε κορυφή
  R1 <- BFS(G, s)                  // ακολουθεί κανονικές κατευθύνσεις
  Εάν R1 ≠ V: επίστρεψε False
  Φτιάξε G^rev = αντιστροφή φοράς όλων των ακμών στο G
  R2 <- BFS(G^rev, s)              // ακολουθεί αντίστροφες
  Εάν R2 ≠ V: επίστρεψε False
  επίστρεψε True
\end{verbatim}
\end{codebox}

\textbf{Διαισθητικά:} το πρώτο BFS ελέγχει αν \textbf{όλες} οι κορυφές είναι
προσπελάσιμες από την $s$· το αντίστροφο BFS (σε ακμές με αντεστραμμένη φορά)
ελέγχει αν η $s$ είναι προσπελάσιμη \textbf{από} κάθε κορυφή.

\textbf{Πολυπλοκότητα:} δύο διασχίσεις BFS σε γραφήματα μεγέθους $O(n + m)$ $\to$
\textbf{$O(n + m)$}.

\begin{intuition}
\textbf{Σημαντικό:} η αντιστροφή του $G$ μπορεί να γίνει σε $O(n + m)$ --- απλά
διασχίζεις όλες τις ακμές και τις προσθέτεις ανάποδα σε νέο γράφημα.

\textbf{Σχετικό:} ο αλγόριθμος του \textbf{Kosaraju} για \textbf{όλες} τις
ισχυρά συνεκτικές συνιστώσες ενός γραφήματος χρησιμοποιεί ακριβώς την ίδια ιδέα
(DFS στο $G$, DFS στο $G^{rev}$).
\end{intuition}

\section{Εισαγωγή στη Συντομότερη Διαδρομή με Βάρη}

Μέχρι τώρα οι ακμές μας ήταν \textbf{χωρίς βάρη} --- κάθε ακμή «μετράει» 1. Με
BFS λύσαμε την απόσταση σε χρόνο $O(n + m)$.

\textbf{Νέα διατύπωση:} σε κάθε ακμή $e$ ανατίθεται \textbf{βάρος}
$\ell_e > 0$ (θετικός πραγματικός). Το \textbf{μήκος μιας διαδρομής} είναι το
άθροισμα των βαρών των ακμών της:
\[
\text{length}(s \to v_1 \to \dots \to t) = \ell_{s,v_1} + \ell_{v_1,v_2} + \dots
+ \ell_{v_{k-1}, t}
\]

\textbf{Πρόβλημα συντομότερης διαδρομής (s--t shortest path):} δεδομένων $s, t$
και βαρών, βρες τη διαδρομή από $s$ σε $t$ με το \textbf{ελάχιστο} συνολικό
βάρος. (Παράδειγμα: σε ένα τυπικό γράφημα, η συντομότερη διαδρομή μπορεί να είναι
$s \to 2 \to 3 \to 5 \to t$ με μήκος $9 + 21 + 2 + 16 = 48$ --- παρότι περνάει
από περισσότερες ακμές από μια εναλλακτική.)

\subsection{Γιατί το BFS δεν αρκεί;}

\textbf{Παρατήρηση:} σε γράφημα με βάρη, η συντομότερη διαδρομή δεν είναι
αναγκαστικά αυτή με τον μικρότερο \textbf{αριθμό ακμών}. Μπορεί να συμφέρει να
πάρω \textbf{3 «φτηνές»} ακμές αντί για \textbf{1 «ακριβή»}.

\textbf{Παράδειγμα:} $s \to t$ με μία ακμή βάρους 100. Εναλλακτικά
$s \to a \to b \to t$ με τρεις ακμές βάρους 10 η καθεμία (σύνολο 30). Το BFS θα
προτιμούσε την πρώτη (1 ακμή), αλλά είναι \textbf{χειρότερη} από άποψη βάρους.

\begin{keypoint}
\textbf{Χρειαζόμαστε νέους αλγορίθμους:}
\begin{itemize}
  \item \textbf{Dijkstra} --- λύνει shortest paths από μία πηγή σε όλους τους
  προορισμούς, \textbf{αν όλα τα βάρη είναι θετικά}. Χρόνος: $O((n + m) \log n)$
  με heap.
  \item \textbf{Bellman-Ford} --- επιτρέπει και \textbf{αρνητικά} βάρη. Χρόνος:
  $O(n \cdot m)$. Εντοπίζει αρνητικούς κύκλους.
\end{itemize}
Θα τους δούμε στο επόμενο μάθημα.
\end{keypoint}

\begin{recapbox}
\textbf{Τι κρατάμε από αυτή τη διάλεξη}
\begin{enumerate}
  \item \textbf{Εφαρμογές BFS/DFS:} συνεκτικές συνιστώσες, flood fill,
  διμερότητα --- όλες σε $O(n+m)$.
  \item \textbf{Διμερότητα $\leftrightarrow$ απουσία περιττού κύκλου.} Έλεγχος:
  BFS $+$ αν ακμή ενώνει ίδιο επίπεδο $\to$ όχι διμερές.
  \item \textbf{Κατευθυνόμενα γραφήματα:} ακμές με κατεύθυνση, εσώβαθμος /
  εξώβαθμος. BFS/DFS επεκτείνεται φυσικά. Web crawler.
  \item \textbf{Ισχυρή συνεκτικότητα:} δύο BFS --- ένα στο $G$, ένα στο
  $G^{rev}$, από οποιαδήποτε $s$. $O(n+m)$.
  \item \textbf{Shortest path με βάρη:} προετοιμασία για Dijkstra /
  Bellman-Ford. Το BFS \textbf{δεν} αρκεί.
\end{enumerate}
\end{recapbox}

\lecturefooter
\end{document}
